% =====================================================================
%  Starter post: copy this file, pick a theme, write.
%  Every command below works the same in all eleven themes.
% =====================================================================
\documentclass[12pt]{beamer}

% Pick one:
%   igsignature   igpreprint    igchalkboard  ignocturne
%   igblueprint   ignotebook    igposter      igtextbook
%   igsorbet      igarchive     igflashcard
\usetheme{igsignature}

\ighandle{@quantum_danesh}      % shown on every slide
\igseries{Qubits desde cero}    % series name ...
\igissue{05}                    % ... and post number
% \iglang{en}                   % English UI words (swipe, Save, ...)

\begin{document}

% Optional extra text some covers show (abstract, chapter summary, ...)
% \coverblurb{...}
\coverslide{Kicker}{Título del post}{Subtítulo de una línea.}

% {Title}{Kicker}: the kicker is the small label above the title
\begin{frame}{La regla de Born}{Ejemplo}
Si medimos $\ket{\psi} = \alpha\ket{0} + \beta\ket{1}$ en la base
computacional, obtenemos $0$ o $1$ con probabilidades
\[
  P(0) = \explain{|\alpha|^2}{peso de $\ket{0}$},
  \qquad
  P(1) = \mhl{|\beta|^2}.
\]
Como el estado está normalizado, \hl{las probabilidades suman uno}:
\keyeq{|\alpha|^2 + |\beta|^2 = 1}
\end{frame}

% Boxes: defbox, thmbox, notebox, eqbox -- all take an optional [title]
\begin{frame}{Cajas y listas}{Componentes}
\begin{defbox}[Definición]
Un \key{qubit} es un sistema cuántico de dimensión $2$.
\end{defbox}
\begin{itemize}
  \item Una viñeta con $\Tr\rho = 1$.
  \item Otra viñeta, con \hl{resaltado}.
\end{itemize}
\begin{notebox}[Nota]
Las notas sirven para matices y advertencias.
\end{notebox}
\end{frame}

% Overlays: every step becomes its own carousel slide
\begin{frame}{Paso a paso}{Animación}
\begin{align*}
  \rho &= \ketbra{\psi}\\
       &\uncover<2->{= |\alpha|^2\ketbra{0} + \alpha\beta^*\ket{0}\!\bra{1} + \dots}
\end{align*}
\uncover<3->{Los términos cruzados son las \hl<3->{coherencias}.}
\end{frame}

\closingslide{Gracias por leer}
  {Una frase final o lo que viene en el próximo post.}

\end{document}
